\subsection{Series estimates for \texorpdfstring{$1<\lambda<2/\sqrt3$}{1 < lambda < 2/sqrt(3)}}
\label{subsec:low-computer-region}
Fix $1<\lambda<\min(m,2/\sqrt3)$. Lemma~\ref{lem:calibration}
gives $D_\lambda>0$ when $\alpha=1$. For $0<\alpha<1$, put
\[\begin{gathered}
 q=\frac{u}{1-u}=\frac hc,\qquad R_* =\frac vc,\qquad
 w_i=\lambda^{i-1}\beta_i^2,\\
 T_p=\sum_{i\ge1}w_i,\qquad Z_p=2q-(\lambda+1)T_p,\qquad
 R=\frac{T_p}{Z_p},\quad \theta=\frac Rq.
\end{gathered}\]
The denominator $Z_p$ is positive by
Lemma~\ref{lem:harmonic-speed}, which also gives $R_*<R<q$.
Write
\begin{equation}\label{eq:low-reference-calibration}
 \alpha_p=\frac{u-R\{1-(\lambda+1)u+\lambda\beta_2\}}{\beta_2}
 =\frac{\lambda^2q+(1-\theta)d(\lambda+1)}{q+1},
 \qquad C_\alpha=\alpha_p+\lambda R.
\end{equation}
Thus $\lambda^2u<\alpha_p<a(cR)<\alpha<1$.  In particular,
$0<h<1/(\lambda+1)<1/2$. Throughout the comparison, the offspring
law is fixed.

\subsubsection{A lower bound for \texorpdfstring{$F_{\rm floor}(v)$}{the moment series at v}}
We bound the series for $F_{\rm floor}(v)$ in terms of
$e=1-\rho_2$. We keep the term at $i=1$, the boundary terms from
summation by parts, and the contribution of $R-R_*$.
For $i\ge1$ define
\begin{align}
 L_i&=\frac1\lambda-r_{i-1}
       =\frac{cd}{\lambda(\lambda^i-d)},\nonumber\\
 W_i^0&=\frac{ci}{\lambda}-1+2iL_i,\nonumber\\
 W_i^p&=cW_i^0+4iL_i
       -2r_i\{iL_i+(i+1)\lambda L_{i+1}\},
 \nonumber\\
 W_i(R)&=W_i^0+RW_i^p,\label{eq:low-W}\\
 V_i&=\sum_{s>i}w_sW_s(R),
 \nonumber\\
 F_{\rm pipe}&=-2RuL_1+u^2W_1(R)+V_1,\nonumber\\
 \tau_i&=1-C_\alpha r_i+(\lambda+1)R-R/r_{i-1},\nonumber\\
 B_i&=w_i\left\{\frac{2i(\alpha_p+cR)}\lambda
       +2iRr_iL_i+2(i+1)R\lambda r_i^2L_i
       +k(2i-1)\tau_i\right\},
 \nonumber\\
 S_i&=B_i-2V_i.\label{eq:low-B-S}
\end{align}
Set $W_1^Y=0$, $W_i^Y=k w_i(2i-1)R/r_{i-1}$ for $i\ge2$, and
\begin{equation}\label{eq:low-J}
 J_i=V_i+k w_i(2i-1)C_\alpha r_i
                         -\frac{W_i^Y+W_{i+1}^Y}{2},\qquad i\ge2.
\end{equation}
Here $J_i$ is the coefficient of $\Delta_i^2$ in the bound below.
The polynomial $J_{s-1}(w)$ in~\eqref{eq:assigned-J} carries the
argument $w$.

\begin{lemma}[A bound in terms of the differences $\Delta_i$]\label{lem:calibrated-head}
Suppose $0<\alpha<1$. Then $V_i>0$ for every
$i\ge1$, $\tau_1=0$, and $0<\tau_i<1$ for $i\ge2$.  Put
\[
 \kappa=\frac{(1-u)T_p}{u^2},
\]
\[\begin{split}
 C_{\rm head}={}&2A_0u\{-k\log(a_2/u)-1+a_2/u\}\\
 &+\frac{(u-\lambda a_2)
       \{2-(3\lambda+1)u+\lambda(\lambda+1)a_2\}}\lambda.
\end{split}\]
Then $C_{\rm head}>0$, $\kappa>1$, and
\begin{align}
 F_{\rm floor}(v)>{}&F_{\rm pipe}
 +\sum_{i\ge1}\Delta_i\rho_i S_i
 +e^2\left\{V_1+\frac{2u^2 C_\alpha}{\lambda(1-e)}
                         -\frac{W_2^Y}{2}\right\}
 +\sum_{i\ge2}\Delta_i^2J_i+C_{\rm cal}e,
 \label{eq:low-master}\\
 C_{\rm cal}&=\frac{C_{\rm head}u^2(\kappa-1)(2-e)}{Z_p}>0.
 \nonumber
\end{align}
The coefficient in braces in~\eqref{eq:low-master} is positive.  Also
\begin{equation}\label{eq:low-e-caps}
 \begin{gathered}
 0<e< E_{\rm hard}=\frac{1-\alpha_p}{1+\lambda R},\qquad
 e\le E_{\rm var}=\frac{d}{(q+1)^2},\\
 e q^2<\frac{\{1-(\lambda+1)h\}u^2}{1+cu},\qquad
 \sqrt e<\frac{d}{q+1}.
 \end{gathered}
\end{equation}
Every series in~\eqref{eq:low-master} converges absolutely.
\end{lemma}
\begin{proof}
First compare $F_{\rm floor}(v)$ with $F_{\rm floor}(cR)$.
With the moments fixed, differentiate $t_n(cR)$ in~\eqref{eq:assigned-t}
with respect to $R$. If the law of $U$ is weighted by $U^2/a_2$, then
\[\begin{aligned}
 \frac{\dd}{\dd R}t_n(cR)
 &=\frac{A_0}{a_2}a_{n+2}
       -\E[U^n(1-U)(1-\lambda U)]\\
 &=-a_2\operatorname{Cov}_{U^2/a_2}
       \left(U^n,\frac{(1-U)(1-\lambda U)}{U^2}\right)>0
\end{aligned}\]
for $n\ge1$.  The second function is strictly decreasing on
$(0,1/\lambda)$; all its weighted expectations are finite.  The
derivative is also bounded above by $(A_0/a_2)a_{n+2}$.
Thus the difference $\widetilde F(cR)-F_{\rm floor}(cR)$ in~\eqref{eq:strict-Mellin-gap} increases with
$R$, and the geometric bounds in Lemma~\ref{lem:mellin-floor}
justify termwise differentiation.

The first two terms in $-c\widetilde F'(w)$ have the lower bounds
\begin{align*}
 2c^2\E[g\cJ(1-y)]&\ge2A_0\gamma_0,\\
 c\E[b\cJ((1-y)(2-y))]&\ge
 \frac{(u-\lambda a_2)
       \{2-(3\lambda+1)u+\lambda(\lambda+1)a_2\}}\lambda.
\end{align*}
They follow by the $U,V$ substitution and
$(1-\lambda UV)^{-2}\ge1$; the remaining factors are positive.
Jensen's inequality under the probability weighted by $U/u$ gives
$\gamma_0\ge u\{-k\log(a_2/u)-1+a_2/u\}$.
The expression in braces is positive: the function
$-k\log z-1+z$ vanishes at $z=1/\lambda$ and has negative derivative
for $0<z<1/\lambda<k$. The term involving $\Gamma$ in
$-c\widetilde F'(w)$ makes the bound strict. Consequently
\begin{equation}\label{eq:low-slope-credit}
 F_{\rm floor}(v)>F_{\rm floor}(cR)+C_{\rm head}(R-R_*).
\end{equation}

We next bound $R-R_*$ for the same offspring law. The identity
$\lambda^{s-1}\beta_s\beta_{s-1}=q(\beta_{s-1}-\beta_s)$,
finite summation by parts, and the vanishing last endpoint give
\begin{align*}
 q-Q&=q\sum_{i\ge1}\Delta_i(a_i+a_{i+1}),\\
 T_p-T&=\sum_{i\ge1}\Delta_i(\rho_i+\rho_{i+1})
                              \sum_{s>i}w_s,\\
 R-R_*&=\frac{2q}{Z_p Z}\sum_{i\ge1}\Delta_i
 \left\{(\rho_i+\rho_{i+1})\sum_{s>i}w_s-T_p(a_i+a_{i+1})\right\}.
\end{align*}
All tail terms in the first two formulas are nonnegative.  The ratios
$r_i$ increase strictly.  With $f_i=(\sum_{s>i}w_s)/\beta_i$, the same
telescoping sum gives
\[
 f_i-f_{i+1}=\lambda^i
 \left\{\frac{\beta_{i+1}^2}{\beta_i}
               -\frac1q\sum_{s>i+1}w_s\right\}<0.
\]
Indeed, each term of the last tail is strictly larger than $r_i$
times the corresponding term $\lambda^{s-1}\beta_s\beta_{s-1}$.
Summing the latter over $s>i+1$ gives $q\beta_{i+1}$. Therefore
$f_i\ge f_1=T_p/u-u=T_p+u(\kappa-1)$.
Solving $\alpha_p<1$ in terms of $\kappa$ gives
$\kappa>2\lambda/(\lambda+1)>1$.  It follows that every expression
in braces in the speed-difference formula is strictly larger than
$u(\kappa-1)\beta_i(\rho_i+\rho_{i+1})$.
Since $0<Z<2q$ and
\[
 \sum_{i\ge1}\beta_i(\rho_i^2-\rho_{i+1}^2)
 \ge u(1-\rho_2^2)=ue(2-e),
\]
we conclude that
\begin{equation}\label{eq:low-speed-deficit}
 R-R_*>\frac{u^2(\kappa-1)e(2-e)}{Z_p}.
\end{equation}

To identify the coefficients of $\Delta_i^2$, put
$\bar\alpha=a(cR)$. The definitions of $a(cR)$ and $\alpha_p$ give
\[
 \bar\alpha-\alpha_p=\frac{C_\alpha e}{1-e}.
\]
For $n_i=a_{i+1}/a_i$ define
$c_i=2i/\lambda-k(2i-1)n_i>1/\lambda$.
Substitute $a_i=\beta_i\rho_i$ and
$\delta_i=\beta_{i+1}\Delta_i$ in the series
\eqref{eq:complete-mellin-floor}.  Use
$(1+L_i)r_i=1/\lambda$ and sum the diagonal terms by parts.  The
result, with every end term displayed, is
\begin{align}
 F_{\rm floor}(cR)={}&F_{\rm pipe}
 +\sum_{i\ge1}\Delta_i\rho_i S_i+\sum_{i\ge1}\Delta_i^2V_i
 +(\bar\alpha-\alpha_p)\sum_{i\ge1}w_i\rho_i\Delta_i c_i
 \nonumber\\[-2pt]
 &+k\sum_{i\ge1}w_i(2i-1)
 \left\{C_\alpha r_i\Delta_i^2
                  -\frac{R\Delta_i\Delta_{i-1}}{r_{i-1}}\right\},
 \qquad\Delta_0=0.\label{eq:low-square-identity}
\end{align}
One may check the local expansion before summation by noting that its
coefficient of $\Delta_i$ differs from $\rho_i B_i$ by
\[
 w_i\rho_i(\bar\alpha-\alpha_p)c_i
 +k w_i(2i-1)\left\{C_\alpha r_i\Delta_i
                         -R\Delta_{i-1}/r_{i-1}\right\}.
\]
More explicitly, before replacing the diagonal terms by the $V_i$
tails, the difference between the $i$th summand of $F_{\rm floor}(cR)$ and these local
terms is
\begin{align*}
 &-(i-1)(1+cR)w_i\rho_i^2
       +i(1+cR)w_{i+1}\rho_{i+1}^2\\
 &\hspace{5mm}
 -\frac{2iRw_iL_i}{r_{i-1}}\rho_{i-1}\rho_i
       +2(i+1)Rw_{i+1}L_i\rho_i\rho_{i+1}.
\end{align*}
The first pair telescopes to zero.  The second pair telescopes because
$L_{i+1}/r_i=L_i$, leaving exactly
$-2Rw_1L_1/r_0=-2RuL_1$.  Its upper endpoint tends to zero
geometrically.  Finally the diagonal telescope uses
$1-\rho_s^2=\sum_{i<s}\Delta_i(\rho_i+\rho_{i+1})$.
This verifies~\eqref{eq:low-square-identity}, including the boundary
term at $i=1$ and the limit at infinity.

At $i=1$ the two $k$ contributions cancel exactly, leaving
$e^2\{V_1+2u^2C_\alpha/[\lambda(1-e)]\}$.
For each $i\ge2$ apply
$2\Delta_i\Delta_{i-1}\le\Delta_i^2+\Delta_{i-1}^2$ to the
remaining adjacent product. The term at $i=2$ contributes
$-W_2^Y e^2/2$, and the other squares have
the coefficients~\eqref{eq:low-J}.  The other terms involving
$\bar\alpha-\alpha_p$ are positive and may be dropped.
Equations~\eqref{eq:low-slope-credit} and
\eqref{eq:low-speed-deficit} now prove~\eqref{eq:low-master}.
The tail positivity $V_i>0$ is proved in the next lemma.  Also
$C_\alpha/(\lambda R)>1$ and $k\lambda r_1<1$, so
\[
 \frac{2u^2C_\alpha}{\lambda(1-e)}>2u^2R,
 \qquad \frac{W_2^Y}{2}<\frac32u^2R.
\]
This proves positivity of the coefficient in braces in~\eqref{eq:low-master}.
Under the distribution $\PP(U_p=\lambda^{-j})=hd^{j-1}$, the quadratic
$Q_{\alpha_p}$ has zero mean and changes sign once from negative
to positive. The proof of Lemma~\ref{lem:mellin-floor} applied to this
scalar distribution gives $\tau_1=0$ and $\tau_i>0$ for $i\ge2$.
Furthermore
\[
 (\lambda+1)-\lambda r_i-1/r_{i-1}
 =-\frac{(\lambda+1)L_i(c+\lambda L_i)}
          {(1+L_i)(1-\lambda L_i)}<0,
\]
so $\tau_i<1$.

Substituting $\bar\alpha-\alpha_p=C_\alpha e/(1-e)$ into $\bar\alpha<1$
gives $e<E_{\rm hard}$. Jensen's inequality $a_2\ge u^2$ gives
$e\le1-u^2/\beta_2=(1-u)(1-\lambda u)=d/(q+1)^2$.
Writing
\[
 E_{\rm hard}=\frac{\theta d(\lambda+1)-\lambda}
                   {(q+1)(1+\lambda\theta q)}
\]
shows that it increases with $\theta$.  Use $\theta<1$ to obtain
the third bound in~\eqref{eq:low-e-caps}.
Finally, $d^2-\{1-(\lambda+1)h\}=h^2+ch>0$ gives the stated
bound on $\sqrt e$.  Absolute convergence follows from
$\beta_i,a_i\le\lambda^{-i}$ and the polynomial factors in $i$;
finite summation followed by geometric domination justifies all
signed rearrangements in this proof.
\end{proof}

\subsubsection{\texorpdfstring{Bounds for the tails $V_i$ and $\sum_{s>i}w_s$}{Bounds for V(i) and the tail sum of w(s)}}
We record the estimates that turn the infinite sums into finite rational
expressions.  Let $\delta=\log\lambda$ and $a=c/\delta$.  Taylor's
formula with integral remainder applied to
$a=\int_0^1(1+c)^s\dd s$ gives
\begin{equation}\label{eq:low-log-enclosure}
 a_L=1+c/2-c^2/12\le a\le
 a_U=1+c/2-c^2/12+c^3/24\le1+c/2.
\end{equation}
For every $z>0$,
\begin{align}
 1+z/2+z^2/12-z^4/720
 &\le\frac{z}{1-e^{-z}}\le1+z/2+z^2/12,\label{eq:low-global-F}\\
 1-z^2/12&\le\frac{z^2}{2(\cosh z-1)}
                       \le1-z^2/12+z^4/240.\label{eq:low-global-C}
\end{align}
For example, multiply the two inequalities in~\eqref{eq:low-global-F}
by the positive exponential denominator.  Beyond the vanishing low
powers, their coefficients, with the indicated sign, are respectively
$(n-3)(n-5)(n-6)(n+8)/(720n!)$ for $n\ge7$ and
$(n-3)(n-4)/(12n!)$ for $n\ge5$.  For~\eqref{eq:low-global-C},
the corresponding even coefficients are
$\{n(n-1)-12\}/(6n!)$ for $n\ge6$ and
$(n-6)(n^3-9n-40)/(120n!)$ for $n\ge8$.
All omitted lower coefficients vanish.  Absolute convergence of the
exponential series proves the inequalities on the entire positive axis.

\begin{lemma}[Bounds for the tails of the Lambert series]\label{lem:low-Lambert-tails}
For every $i\ge1$, $V_i>0$.  If $b=\lambda L_i$, put
\[
 r=\frac1{\lambda+b}=r_i,\qquad
 A=c+br,\quad B=br,\quad t=\frac r{b+c},
\]
and define the five polynomials
\begin{align*}
 P_1&=A,\qquad P_2=A(A+B),\qquad P_3=A(A^2+4AB+B^2),\\
 P_4&=A(A^3+11A^2B+11AB^2+B^3),\\
 P_5&=A(A^4+26A^3B+66A^2B^2+26AB^3+B^4).
\end{align*}
Define the following explicit rational functions of an argument $a>0$:
\begin{align}
 T_L(a)&=t\{a+A/2+P_2/(12a)-P_4/(720a^3)\},\nonumber\\
 T_U(a)&=t\{a+A/2+P_2/(12a)\},\nonumber\\
 P_U(a)&=\frac t\lambda
 \{(i+1)(aA+P_2/2+P_3/(12a))+a^2-P_2/12+P_4/(240a^2)\},\nonumber\\
 P_L(a)&=\frac t\lambda
 \{(i+1)(aA+P_2/2+P_3/(12a)-P_5/(720a^3))
                          +a^2-P_2/12\}.\label{eq:low-tail-rationals}
\end{align}
At $a=c/\log\lambda$,
\begin{align}
 V_i/w_i&\le V_U(a)
 :=-(1+cR)T_L(a)+\{1+(\lambda+1)R\}P_U(a)
                   -2R(i+1)\lambda r^2-2Rt,\nonumber\\
 V_i/w_i&\ge-(1+cR)T_U(a)+\{1+(\lambda+1)R\}P_L(a)
                   -2R(i+1)\lambda r^2-2Rt.\label{eq:low-V-bound}
\end{align}
For $q\le10$, $c\le1/6$, and every integer $i\ge1$, these imply the rational
bounds
\begin{equation}\label{eq:low-V-rational-bound}
 V_L\le V_i/w_i\le V_U,
\end{equation}
where
\[\begin{aligned}
 V_U&=V_U(a_U),\\
 V_L&=-(1+cR)T_U(a_U)+\{1+(\lambda+1)R\}P_L(a_L)\\
 &\quad-2R(i+1)\lambda r^2-2Rt.
\end{aligned}\]
The tail $\sum_{s>i}w_s$ lies between $w_iT_L(a_L)$ and
$w_iT_U(a_U)$ on this domain.
\end{lemma}
\begin{proof}
Geometric expansion of $\beta_s^2$ and summation over every $s>i$ give
\[
 \sum_{s>i}w_s=\frac{h^2}{\lambda\delta}
 \sum_{N\ge1}d^{N-1}\lambda^{-N(i+1)}
                         \frac{N\delta}{1-e^{-N\delta}}.
\]
Hold $h,d$ fixed while differentiating this deterministic series with
respect to $\lambda$.  A direct derivative of its $N$th summand gives
the identity
\begin{align}
 V_i={}&-(1+cR)\sum_{s>i}w_s-2R(i+1)w_{i+1}-\frac{2Rq\beta_{i+1}}\lambda
 \nonumber\\
 &+\{1+(\lambda+1)R\}\frac{ch^2}{\lambda^2}
 \Biggl\{\frac{i+1}{\delta}\sum_{N\ge1}N d^{N-1}\lambda^{-N(i+1)}
                                      \frac{N\delta}{1-e^{-N\delta}}
 \nonumber\\[-2pt]
 &\hspace{39mm}+\frac1{\delta^2}
 \sum_{N\ge1}d^{N-1}\lambda^{-N(i+1)}
                          \frac{(N\delta)^2}{2(\cosh(N\delta)-1)}\Biggr\}.
 \label{eq:low-full-tail-identity}
\end{align}
The terms $-2R(i+1)w_{i+1}$ and $-2Rq\beta_{i+1}/\lambda$ come from
shifting the summation indices in the terms involving $W_i^p$.
For a direct check of the differentiation, write
$y=\lambda^{-N}$ and
$T_N=h^2N d^{N-1}y^{i+1}/[\lambda(1-y)]$.  Then
\[
 T_N'=-T_N/\lambda-
 \frac{h^2N^2d^{N-1}y^{i+1}}{\lambda^2}
 \left\{\frac{i+1}{1-y}+\frac{y}{(1-y)^2}\right\}.
\]
Substitution proves~\eqref{eq:low-full-tail-identity} term by term.
On a compact neighborhood of a fixed $\lambda>1$, the summands and
their derivatives are bounded by a geometric sequence times a fixed
polynomial in $N$.  This justifies the derivative and both summations.

For $p=\lambda^{i+1}$, the geometric moments are
$M_j=\sum_{N\ge1}N^j d^{N-1}p^{-N}$, where
$M_0=(p-d)^{-1}$ and $M_{j+1}=-p\partial_p M_j$.
Applying this recurrence up to $j=5$ gives
\[
 \frac{hM_0}{w_i}=\frac{\lambda r}{q(b+c)},\qquad
 \frac{h^{j+1}M_j}{w_i}
    =\frac{q^{j-1}\lambda rP_j}{b+c}\quad(1\le j\le5).
\]
Use~\eqref{eq:low-global-F}--\eqref{eq:low-global-C} in every
$N$ summand of~\eqref{eq:low-full-tail-identity}.  The sign of the
tail coefficient is negative and those of the two derivative kernels
are positive.  The resulting inequalities are precisely
\eqref{eq:low-tail-rationals}--\eqref{eq:low-V-bound}.

For $q\le10$ we have $c\le1/6$ and
$b\le d/(q+i)$.  Hence $A\le2/3$, $0\le B\le A$,
$P_2\le8/9$, $P_3/A\le8/3$, and $P_4\le128/27$.
Differentiation of the explicit rational functions gives
\[
 0<T_L'(a)\le\frac{11t}{10},\qquad
 P_U'(a)\ge\frac{9t}{5\lambda},\qquad
 T_U'(a)>0,\quad P_L'(a)>0\qquad(a\ge1).
\]
For instance, $T_L'/t=1-P_2/(12a^2)+P_4/(240a^4)$,
and $\lambda P_U'/t=(i+1)(A-P_3/(12a^2))+2a-P_4/(120a^3)$;
the displayed polynomial bounds prove these estimates.
The derivative of $V_U(a)$ is positive since
\[
 \frac95\{1+(\lambda+1)R\}-\frac{11}{10}\lambda(1+cR)>0
 \qquad(1<\lambda\le7/6, R\ge0).
\]
Now~\eqref{eq:low-log-enclosure} proves
\eqref{eq:low-V-rational-bound} and the tail enclosure.

To prove positivity, put $z=d\lambda^{-i}$ and
$\Pi(N)=N(1-z)^2z^{N-1}$ for $N\ge1$.  This is a probability
distribution.  The same geometric summation, including the two
endpoints, gives
\begin{equation}\label{eq:low-positive-mixture}
 \frac{V_i}{w_i}=\frac1\lambda\E_\Pi
       [i(A_N+RP_N)+B_N+RC_N],
\end{equation}
where, writing $Q_N=1+\lambda+\cdots+\lambda^{N-1}$,
\[\begin{gathered}
 A_N=\frac{N}{Q_N},\quad B_N=\frac{\lambda Q_N'}{Q_N^2},\quad
 P_N=(\lambda+1)A_N-2\lambda^{1-N},\\
 C_N=(\lambda+1)B_N-
          \frac{2\lambda(1-A_N)}{cN\lambda^N}.
\end{gathered}\]
Here $A_N>0$, $B_N\ge0$, and
$P_N\ge cN\lambda^{-N}>0$: the latter is the convex-chord
inequality $N(\lambda^{N+1}+1)\ge2\sum_{j=1}^N\lambda^j$.
Also $Q_N-N\le cQ_N'$ term by term, and hence
\[
 C_N\ge B_N\left\{(\lambda+1)-\frac{2Q_N}{N\lambda^N}\right\}\ge0.
\]
This proves $V_i>0$ and completes the proof.
\end{proof}

The tail estimates also provide power-sum enclosures used below.  For
$A=q+2$ and every integer $j\ge2$, put
\begin{equation}\label{eq:low-power-enclosures}
 U_j=\frac{A^{1-j}}{j-1}+\frac1{2A^j}+\frac{j}{12A^{j+1}},\qquad
 L_j=U_j-\frac{j(j+1)(j+2)}{720A^{j+3}}.
\end{equation}
Then $L_j\le\sum_{i\ge2}(q+i)^{-j}\le U_j$.
Indeed, write the sum as
\[
 \frac1{(j-1)!}\int_0^\infty
       \frac{z^{j-1}e^{-Az}}{1-e^{-z}}\dd z
\]
and insert the global bounds~\eqref{eq:low-global-F}; Tonelli and
the gamma integral evaluate every term.  Thus~\eqref{eq:low-power-enclosures}
includes the entire infinite tail.

\subsubsection{The range \texorpdfstring{$0<q\le1/15$}{0 < q <= 1/15}}
Here $u\le1/16$. For the logarithmic term we retain the first two
terms in the definition of $\overline\gamma_{s-1}$. Define
$F_{\rm cap}$ by the right side of~\eqref{eq:complete-mellin-floor},
with $w=v$ and with $\overline\gamma_{s-1}$ replaced by
$a_s(r_s-r_{s-1})-\delta_s$.
The exact logarithmic identity in Lemma~\ref{lem:mellin-floor} gives
$D_\lambda\ge F_*\ge F_{\rm cap}$.
For this paragraph only put
\begin{align*}
 B_s&=\frac{c}{\lambda^s-1},& f_s&=\lambda^{s-1}B_s^2,&
 R_s&=1-u+uB_s,& \frac{\beta_s}{u}&=\frac{B_s}{R_s},\\
 t&=T/u^2,&q_a&=Q/u,&t_\pi&=\sum_{s\ge1}f_s/R_s^2,\\
 H&=\sum_{s\ge1}s f_s^2\rho_s^2/R_s^3,&
 H_\pi&=\sum_{s\ge1}s f_s^2/R_s^3,\\
 p_s&=\frac{cs}{\lambda}\frac{\lambda^s+1}{\lambda^s-1}
                 -\frac{\lambda+1}\lambda,&
 P&=\sum_{s\ge1}f_sp_s\rho_s^2/R_s^2,&
 M&=\frac{2q_a}{2q_a-(\lambda+1)ut}.
\end{align*}
The function $x\coth x$ increases on $(0,\infty)$, because its
derivative has numerator $\sinh(2x)-2x>0$.  Hence $p_s\ge0$.
Finally set
\[
 G_s=\frac{s\lambda^{s-1}(1+L_s)a_s\delta_s}{u^2},\quad
 G=\sum_{s\ge1}G_s,\quad
 G_+=\sum_{s\ge1}G_s\left\{\lambda+
 \frac{(s+1)\lambda a_{s+1}}{s a_s(1+L_s)}\right\}\ge\lambda G.
\]

Separating the term at $s=1$ gives
\begin{equation}\label{eq:low-small-raw}
 \frac{F_{\rm cap}}{u^2}
 =M\left\{P+u\left(\frac{t^2}{q_a}-2H\right)
                 +\frac{ut}{q_a}G_+\right\}+2\alpha G.
\end{equation}
Here is a verification that also specifies the endpoint in this
identity.  Put $H_s^{\rm loc}=s-(s+1)\lambda L_{s+1}$ and
$D_s^{\rm loc}=2sL_s-s/\lambda+H_s^{\rm loc}r_s$. Expanding
$F_{\rm cap}$ gives
\begin{align*}
 F_{\rm cap}={}&-2R_*uL_1
 +\sum_{s\ge1}\lambda^{s-1}a_s^2
             \{(1+v)W_s^0+2R_*D_s^{\rm loc}\}\\
 &+2\sum_{s\ge1}s\lambda^{s-1}a_s\delta_s
 \left\{\alpha(1+L_s)+R_*\left(c+\lambda L_s
                         +\frac{(s+1)\lambda L_{s+1}}s\right)\right\}.
\end{align*}
The elementary identities
\[
 W_s^0-1/\lambda=p_s-2us f_s/R_s
\]
and, writing $t_N=u^{-2}\sum_{s=1}^N\lambda^{s-1}a_s^2$ and
$Q_N=\sum_{s=1}^N\lambda^{s-1}a_sa_{s-1}$,
\begin{align*}
 \sum_{s=1}^N\frac{\lambda^{s-1}a_s^2}{u^2}
                  (D_s^{\rm loc}-W_s^0)
 ={}&t_N-1-\frac{Q_{N+1}-u}{\lambda u^2}\\
 &+\frac1{u^2}\sum_{s=1}^N\lambda^{s-1}a_s\delta_s
           \left\{H_s^{\rm loc}+(s+1)\lambda\frac{a_{s+1}}{a_s}\right\}\\
 &+\frac{(N+1)\lambda^Na_{N+1}^2}{u^2}
\end{align*}
give~\eqref{eq:low-small-raw} after substitution of
$R_*=ut/[2q_a-(\lambda+1)ut]$.  The last endpoint tends to zero
geometrically. The middle sum involving $\delta_s$ combines with
the other terms proportional to $R_*$ to give the term containing $G_+$.

We next bound the loss $t_\pi-t$ by $G$.  With
$g_i=i\lambda^{i-2}\beta_i^2$,
$u^2G=\sum_i g_i\rho_i\Delta_i$. For the reference moments $\beta_i$,
\[
 \sum_{n>i}\lambda^{n-1}\beta_n\beta_{n-1}=q\beta_i,
 \qquad \sum_{n>i}w_n<q\beta_i/\lambda,
 \qquad \frac{q\beta_i}{g_i}
 =\frac{\lambda^2(1-d\lambda^{-i})}{ci}
 \le\frac{\lambda^2}{1-u}.
\]
The last bound follows by Bernoulli's inequality and
$h=cu/(1-u)$. Summing $1-\rho_s^2=\sum_{i<s}\Delta_i(\rho_i+\rho_{i+1})$
against $w_s/u^2$ and using the tail bounds gives
\begin{equation}\label{eq:low-small-relative-loss}
 t_\pi-t\le\frac{2\lambda G}{1-u}.
\end{equation}
Also $1\le q_a\le(1-u)^{-1}$ and $t\ge1$.
Writing $f_s=[\sinh(\log\lambda/2)/\sinh(s\log\lambda/2)]^2$
shows that $f_s\le s^{-2}$.  Five terms and an integral give
$\sum s^{-2}<5989/3600<5/3$.  Since $f_1=R_1=1$,
\[
 t_\pi\le1+\frac{2}{3(1-u)^2}<\frac{16}9,\qquad
 \kappa=(1-u)t_\pi\le\frac53,\qquad
 M\le\frac{90}{79},\qquad t/q_a\ge\frac{15}{16}.
\]
The comparison $R_*<R$ and $a_2<\beta_2$ gives
\[
 \alpha>\alpha_p
 =\frac{(\lambda+1)(2-\kappa)-2\lambda u}
              {2-(\lambda+1)u\kappa}
 \ge\frac{2/3-(12/5)u}{2-2u(1-u)}\ge\frac{992}{3615}
 \quad(\lambda\le6/5, u\le1/16).
\]
For the middle inequality the numerator is positive, $\kappa\le5/3$,
and $\kappa\ge1-u$; the last rational function decreases on the
stated interval, since its derivative numerator is a positive constant
times $18u^2-10u-13<0$.

Put $B=(1-u)t_\pi^2-2H_\pi$.
Equation~\eqref{eq:low-small-relative-loss} gives
$t^2/q_a-2H\ge B-4\lambda t_\pi G$.
Retaining $G_+$ in~\eqref{eq:low-small-raw} now yields
\begin{equation}\label{eq:low-small-budget}
 \frac{D_\lambda}{u^2}\ge M(P+uB)+C_*G,\qquad
 C_* =\frac{1984}{3615}
 -\frac{90}{79}\frac65\frac1{16}
       \left(\frac{64}{9}-\frac{15}{16}\right)
 =\frac{389899}{18277440}>0.
\end{equation}
The $s=2$ term in $P$ and the $s=1$ term in $G$ give
\[
 P\ge\frac{c^2\rho_2^2}{(\lambda+1)^3},\qquad
 G\ge\frac{1-\rho_2}\lambda.
\]
Thus~\eqref{eq:low-small-budget} also holds in the form
\begin{equation}\label{eq:low-small-root-budget}
 \frac{D_\lambda}{u^2}\ge
 M\{P+uB+K_0(1-\rho_2)\},\qquad
 K_0=\frac{389899}{24986880}.
\end{equation}

For clarity, the finite rational input to the remaining scalar estimate
is stated separately.  For $\ell>1$ let $c_\ell=\ell-1$,
$f_s(\ell)=\ell^{s-1}c_\ell^2/(\ell^s-1)^2$,
$B_s(\ell)=c_\ell/(\ell^s-1)$, and put
\begin{align*}
 \tau(\ell)&=\frac{c_\ell^2}
 {\ell(c_\ell-c_\ell^2/2+c_\ell^3/3)(\ell^9-1)}
                         +\frac{f_9(\ell)}2,\\
 E(\ell)&=\frac{c_\ell^4\ell^{-16}}{6\ell^2}
 \left\{\frac{8}{2c_\ell/(\ell+1)}
                     \frac{3-\ell^{-8}}{(1-\ell^{-8})^3}
 +\frac1{(2c_\ell/(\ell+1))^2}
       \left(\frac{2-\ell^{-8}}{(1-\ell^{-8})^2}-\frac12\right)\right\}.
\end{align*}
\begin{lemma}[Six rational bounds used for $0<u\le1/16$]\label{lem:low-Q0-arithmetic}
The following rational inequalities hold:
\begin{align*}
 \sum_{s=1}^8 f_s(6/5)+\tau(6/5)&>1559/1000,\\
 \sum_{s=1}^8 f_s(23/20)+\tau(23/20)&>63/40,\\
 \sum_{s=1}^8 f_s(11/10)+\tau(11/10)&>19/12,\\
 \sum_{s=1}^8 f_s(6/5)(1-1/s)+\tfrac89\tau(6/5)&>9/25,\\
 \sum_{s=1}^8s^{-3}\{1-B_s(6/5)\}+1/128&<2/15,\\
 \sum_{s=1}^8s f_s(23/20)^2+E(23/20)&<239/200.
\end{align*}
All denominators in these six finite expressions are positive.
\end{lemma}

We now estimate the tails starting at $s=9$. Both $f_s$ and
$B_s$ decrease with $\lambda$; the assertion for $f_s$ follows from
the increasing function $x\coth x$, and that for $B_s$ from its finite
geometric denominator.  As a function of real $x>0$, $f(x)$ is
decreasing and convex, with
\[\begin{aligned}
 f'&=-\log\lambda\, f\coth(x\log\lambda/2),\\
 f''&=(\log\lambda)^2f
       \{\coth^2(x\log\lambda/2)+\tfrac12\operatorname{csch}^2(x\log\lambda/2)\}.
\end{aligned}\]
Convex trapezoid comparison and
$\log\lambda\le c-c^2/2+c^3/3$ imply
$\sum_{s\ge9}f_s\ge\tau(\lambda)$.
The function $xf(x)^2$ decreases; its logarithmic derivative is less
than $-3/x$.  Expand its integral tail using the positive series for
$(1-e^{-x\log\lambda})^{-4}$.  The identities
\[
 \frac{\binom{j+3}{3}}{j+2}=\frac{(j+1)(j+3)}6,\qquad
 \frac{\binom{j+3}{3}}{(j+2)^2}
       =\frac{j+2-1/(j+2)}6
\]
sum the two geometric parts.  Keeping the $j=0$ negative subtraction
and using $\log\lambda\ge2c/(\lambda+1)$ gives
$\sum_{s\ge9}s f_s^2\le E(\lambda)$.
Every remaining coefficient is positive, so this substitution has the
stated direction.  Consequently, with
\[
 T_0=\sum f_s,\quad H_0=\sum s f_s^2,\quad
 A=\sum f_s(1-B_s),\quad C_{\rm tail}=\sum s f_s^2(1-B_s),
\]
Lemma~\ref{lem:low-Q0-arithmetic}, $B_s\le1/s$, and
$\sum_{s\ge9}s^{-3}\le1/128$ prove
\[
 A>9/25,\quad C_{\rm tail}<2/15,\quad H_0<241/200,
 \quad H_0<239/200\ \hbox{if }\lambda\ge23/20,
\]
together with the three lower bounds for $T_0$ on the respective
intervals of $\lambda$ specified below. The general bound for $H_0$ is
$\sum_{s=1}^6s^{-3}+\int_6^\infty x^{-3}\dd x=86701/72000<241/200$.

The positive geometric expansion and the convex chord for
$(1-uz)^{-3}$ on $0\le z\le1$ give
\[
 B>(1-u)(L+18u/25)^2
       -2\{\overline H+(2/15)((1-u)^{-3}-1)\}.
\]
Use $(L,\overline H)=(19/12,241/200)$,
$(63/40,241/200)$, and $(1559/1000,239/200)$ on
$(1,11/10]$, $[11/10,23/20]$, and $[23/20,6/5]$, respectively.
The right side decreases in $u\le1/16$; its derivative is
\[
 (L+18u/25)(36/25-L-54u/25)-(4/5)(1-u)^{-4}<0.
\]
At $u=1/16$ the three values are
$973799/51840000$, $-21361/3240000$, and
$-566237/16200000$.
The first is positive.  On the other intervals put
$A_1=10/9261$, $A_2=180/79507$, and let $d_1,d_2$ be the absolute
values of the last two displayed fractions.  Since
$c^2/(\lambda+1)^3$ increases for $1<\lambda\le6/5$,
\[\begin{gathered}
 P\ge A_j\rho_2^2,\qquad K_0>2A_2>2A_1,\\
 A_1-d_1/16=\frac{11873177}{17781120000}>0,\qquad
 A_2-d_2/16=\frac{1636194841}{20608214400000}>0.
\end{gathered}\]
For $0<\rho<1$,
$A_j\rho^2+K_0(1-\rho)\ge A_j$.
Equation~\eqref{eq:low-small-root-budget} proves $D_\lambda>0$.
This covers $q\le1/15$, since $2/\sqrt3<6/5$.

\subsubsection{Rational bounds for \texorpdfstring{$1/15<q\le10$}{1/15 < q <= 10}}
We use the tail bounds in Lemma~\ref{lem:low-Lambert-tails} to
form a finite set of rational sign tests.  Their definitions are given
here so that the arithmetic supplement contains only polynomial
expansion and exact coefficient checks.

For $q,c>0$ set $\lambda=1+c$, $d=1-cq$, $u=q/(q+1)$ and
$\theta_g=1-1/[6(q+1)]$. At $i=1$, use
$b=d/(q+1)$ in~\eqref{eq:low-tail-rationals}, and define
\begin{equation}\label{eq:low-theta-enclosures}
 T_{\rm lo}=u^2\{1+T_L(a_L)\},\quad
 T_{\rm up}=u^2\{1+T_U(a_U)\},\qquad
 \theta_{\rm lo,up}=
 \frac{T_{\rm lo,up}}{q\{2q-(\lambda+1)T_{\rm lo,up}\}}.
\end{equation}
In the rational expressions that follow, $R=q\theta$ and
$\alpha_p$ is given by~\eqref{eq:low-reference-calibration}, with
$\theta$ treated as a variable. For variables $b$ and $i$, put
\begin{align}
 r&=(\lambda+b)^{-1},\qquad r_{\rm prev}=(1-b)/\lambda,\nonumber\\
 \tau&=1-(\alpha_p+\lambda R)r+(\lambda+1)R-R/r_{\rm prev},
 \nonumber\\
 \mathcal B(i,b,\theta)&=\frac{2i(\alpha_p+cR)}\lambda
       +\frac{2iRrb}\lambda+2(i+1)Rr^2b
                              +\frac{a_L}{\lambda}(2i-1)\tau,\nonumber\\
 K_H(i,b,\theta)&=\mathcal B(i,b,\theta)-2V_U,\nonumber\\
 K_J(i,b,\theta)&=V_L+\frac{a_L}{\lambda}(2i-1)\alpha_pr
 -\frac{a_U R}{2\lambda}
       \{(2i-1)(1/r_{\rm prev}-\lambda r)+2\lambda r\},
 \label{eq:low-age-rational-tests}\\
 i_{\rm lo}(b)&=\frac{2(d-qb)}{2b+c(d-qb)}.\nonumber
\end{align}
Here $V_U,V_L$ are exactly~\eqref{eq:low-V-rational-bound}, evaluated
at the free $i,b,\theta$.  Both $K_H$ and $K_J$ are affine in $i$
and in $\theta$ separately.

We also need a lower bound for the terms at $i=1$ and the term
$C_{\rm cal}e$ in~\eqref{eq:low-master}. In the next display all tail
functions are at $i=1$, $b=d/(q+1)$, and
$r=(q+1)/(q+2+c)$.  Put
\begin{align*}
 V_0&=-T_U(a_U)+P_L(a_L),\\
 V_R&=-cT_U(a_U)+(\lambda+1)P_L(a_L)-4\lambda r^2-2r/(b+c),\\
 W_0&=c/\lambda-1+2b/\lambda,\qquad
 W_p=cW_0+4b/\lambda
          -2r\{b/\lambda+2d/(q+2+c)\},\\
 \alpha_0&=\{\lambda^2q+d(\lambda+1)\}/(q+1),\qquad
 \alpha_1=-d(\lambda+1)/(q+1),\\
 B_{10}&=2\alpha_0/\lambda,\qquad
 B_{11}=2(\alpha_1+cq)/\lambda+2qrb/\lambda+4qr^2b,\\
 E_v&=d/(q+1)^2,\qquad
 \varepsilon=u-\lambda\beta_2,\qquad
 \phi=\frac{d-c^2(q+2+c)/6}{(q+2+c)(2q+3+c)},\\
 C_0&=\frac{E_v(2-E_v)}{2q}
                 \{2u\phi+(2-h)\varepsilon/\lambda\}.
\end{align*}
Define, for $0\le e<1$,
\begin{align}
 \mathcal A(e,\theta)={}&(1-e)^2(V_0+q\theta V_R)
 +W_0+q\theta W_p-2\theta d/\lambda+e(B_{10}+\theta B_{11})
 \nonumber\\
 &+e^2\left\{\frac{2\{\alpha_0+\theta(\alpha_1+\lambda q)\}}
                     {\lambda(1-e)}
                         -\frac32a_Uq\theta r\right\}
       +C_0\{\theta(2-h)-1\}e.\label{eq:low-root-rational}
\end{align}
Use $U_j,L_j$ from~\eqref{eq:low-power-enclosures} to define
\begin{align}
 K_4={}&\frac{q^2(q+1)^2}\lambda
       \{U_3-qL_4+c(U_2-2qL_3+q^2U_4)\},\nonumber\\
 K_6={}&\frac{q^4(q+1)^2}\lambda
       \{U_5-qL_6+c(U_4-2qL_5+q^2U_6)\},\nonumber\\
 K_7={}&\frac{q^5(q+1)^2}\lambda
       \{U_6-qL_7+c(U_5-2qL_6+q^2U_7)\},\nonumber\\
 K_{\rm sq}={}&\frac{q^4(q+1)^2}\lambda
       \{U_4-2qL_5+q^2U_6+c(U_3-3qL_4+3q^2U_5-q^3L_6)\}.
 \label{eq:low-rational-costs}
\end{align}
For $s=dz/(q+1)$, $e=s^2$, $0\le z\le1$, define
\begin{align}
 \mathcal G_1(\theta)&=\mathcal A(e,\theta)
                     -\tfrac32e(1-e)(1-s)K_4,\nonumber\\
 \mathcal G_2(\theta)&=\mathcal A(e,\theta)
       -4e(1+s)\left\{\left(1-\frac{qs}{1-s}\right)K_6
                            +\frac{qs}{1-s}K_7\right\}
       -8e^2K_{\rm sq}.\label{eq:low-finite-budgets}
\end{align}
The factor $1-s$ in the denominator of $\mathcal G_2$ is strictly
positive throughout the specified domain.

\begin{lemma}[Rational sign tests for $0<q\le3$ and $3\le q\le10$]
\label{lem:low-Q12-arithmetic}
The following finite assertions hold.
\begin{enumerate}
\item On $0<q\le3$, $0\le c\le1/6$, and also on
$q=1/x$, $c=hx$, $1/10\le x\le1/3$, $0\le h\le1/2$,
\[
 0<T_{\rm lo}\le T_{\rm up}
 <\frac{2q^2(6q+5)}{6(q+1)+(\lambda+1)q(6q+5)}.
\]
The tail expressions at $c=0$ mean their rational continuous
extensions.
\item For $0<q\le3$, $0\le c\le1/6$, $0\le v_0\le1$, and
$i=2+j$, $j\ge0$, put $b=dv_0/(q+i)$, $r=(\lambda+b)^{-1}$ and
\[
 K_R=icr^2+2ibr^3-\frac{(2i-1)\lambda}{2(1-b)}
                         +\frac{(2i-3)\lambda r}{2}.
\]
Then $K_R<0$ and
\[
 i\lambda r^2+(2i-1)\lambda u r+qK_R>0.
\]
For $b=dv_0/(q+2)$ and $\theta\in\{0,\theta_g\}$ one also has
\[
 \partial_iK_H>0,\qquad
 K_H(i_{\rm lo}(b),b,\theta)+\tfrac32(qb/d)^2>0.
\]
The latter rational expressions are used when $c>0,v_0>0$;
their cleared polynomials include the closed boundary.
\item On $q=1/x$, $c=hx$, $1/10\le x\le1/3$,
$0\le h\le1/2$, $b=dv_0/(q+2)$, $0\le v_0\le1$,
at each $\theta\in\{0,\theta_g\}$,
\[
 \partial_iK_H\ge0,\quad K_H(i_{\rm lo},b,\theta)+4(qb/d)^4\ge0,
 \qquad
 \partial_iK_J\ge0,\quad K_J(i_{\rm lo},b,\theta)+8(qb/d)^4\ge0.
\]
\item For $1/29\le q\le3$, $0\le c\le1/6$, $0\le z\le1$,
both $\mathcal G_1(\theta_{\rm lo})$ and
$\mathcal G_1(\theta_{\rm up})$ are strictly positive.
For $q=1/x$, $c=hx$, $1/10\le x\le1/3$,
$0\le h\le1/2$, $0\le z\le1$, both
$\mathcal G_2(\theta_{\rm lo})$ and
$\mathcal G_2(\theta_{\rm up})$ are strictly positive.
\end{enumerate}
\end{lemma}

All the tests are on rational functions explicitly defined above.
Their denominators are products of positive numerical constants and
$\lambda$, $q$, $q+1$, $q+2$, $q+2+c$, $2q+3+c$, $a_L$, $a_U$,
$1-b$, $\lambda+b$, $b+c$, $2b+c(d-qb)$, $1-e$, and $1-s$.
The speed endpoint denominators are positive by part~(i).
After substituting $b=dv_0/(q+2)$, the only additional vanishing
boundary factors are $dv_0+c(q+2)$ and
$(2-cq)v_0+c(q+2)$. Both are strictly positive when
$c>0$ and $v_0>0$, as required for clearing denominators.

For specificity, the finite coefficient operation used to check this
lemma is as follows.  Map each bounded variable to $y\in[0,1]$ by
the affine maps in the lemma.  If
$P(y_1,\ldots,y_d)=\sum_\nu p_\nu y^\nu$ has coordinate degrees
$D_1,\ldots,D_d$, its tensor Bernstein coefficients are
\begin{equation}\label{eq:low-Bernstein-transform}
 b_\kappa=\sum_{\nu\le\kappa}p_\nu
             \prod_{j=1}^d\frac{\binom{\kappa_j}{\nu_j}}
                                  {\binom{D_j}{\nu_j}}.
\end{equation}
The unbounded variable $j=i-2$ in part~(ii) stays in the ordinary
power basis.  Positive coefficients in those variables imply the
claimed sign because their basis functions are nonnegative and sum
to one in each bounded variable.  For the exceptional one-variable
groups, the arithmetic supplement gives the polynomial and its
Euclidean Sturm sequence: positive endpoint values and equal numbers
of endpoint sign variations imply positivity throughout $[0,1]$.
There are fifteen degree-seven groups in part~(iv) for
$\mathcal G_1$, and twenty-five degree-ten groups in part~(iii) for
$K_H$.  Vanishing factors at zero or one are removed explicitly
before applying Sturm's theorem.  The remaining part~(iv) tests have
$10150$ and $10850$ positive tensor coefficients.

We now prove that this finite lemma suffices.  Part~(i) and the
tail bounds for $\sum_{s>1}w_s$ imply
\[
 0<\theta_{\rm lo}\le\theta\le\theta_{\rm up}<\theta_g<1.
\]
For every integer $i\ge2$, Bernoulli's inequality gives
$b=cd/(\lambda^i-d)\le d/(q+i)$.  Also
$\lambda^i=d(1+c/b)$ and
$\log\lambda\le c$, $\log z\ge2(z-1)/(z+1)$ for $z\ge1$
give $i\ge i_{\rm lo}(b)$.  In the definition of $B_i$ its
$k$ coefficient is nonnegative, because $\tau_i>0$.
Hence $S_i/w_i\ge K_H(i,b,\theta)$.
For the square coefficient, expand~\eqref{eq:low-J} as
\[
 J_i/w_i=V_i/w_i+k(2i-1)\alpha_pr
       -\frac{kR}{2}\{(2i-1)(1/r_{\rm prev}-\lambda r)+2\lambda r\}.
\]
The bracket is positive.  Since $a_L/\lambda\le k\le a_U/\lambda$,
$J_i/w_i\ge K_J(i,b,\theta)$. Affinity in $\theta$ and the
nonnegative derivatives with respect to $i$ in the lemma give
\begin{equation}\label{eq:low-all-age-prices}
 \begin{cases}
 S_i\ge-\tfrac32w_i\beta_i^2,\quad J_i>0,&q\le3,\\
 S_i\ge-4w_i\beta_i^4,\quad J_i\ge-8w_i\beta_i^4,&3\le q\le10,
 \end{cases}\qquad i\ge2.
\end{equation}
For $q\le3$, positivity of $J_i$ follows from the lower bound in
\eqref{eq:low-positive-mixture},
\[
 V_i/w_i>i\lambda r_i^2
                      \{1+R(c/\lambda+2L_{i+1})\}.
\]
Indeed $A_N\ge\lambda^{1-N}$ and
$P_N\ge cN\lambda^{-N}$, and summing these two geometric
expectations gives the display.  Use $k\ge1/\lambda$ on the
positive $\alpha_p$ term and $k\le1$ on its negative bracket.
The result is greater than
$i\lambda r^2+(2i-1)\lambda ur+qK_R$, since
$\alpha_p>\lambda^2u$, $R<q$, and $K_R<0$.
Part~(ii) proves positivity.

To bound the negative terms in~\eqref{eq:low-master}, we use
the representation of all the $\Delta_i$ by the same measure $\sigma$. For
$p=2$ or $4$, the sequence
$m_i=w_i\beta_i^p$ is a positive mixture of powers:
\[
 m_i=\frac{h^{p+2}}\lambda
       \sum_{N\ge p+1}\binom{N}{p+1}d^{N-p-1}\lambda^{-Ni}.
\]
For $0\le t<1$, the function
$f_t(z)=z/(1-t+t\sqrt z)^2$ is concave on $[0,1]$, since
\[
 f_t''(z)=-\frac{3(1-t)t}
                       {2\sqrt z(1-t+t\sqrt z)^4}\le0
 \quad(z>0).
\]
In the integral representation of $\Delta_i$ take $\xi$ to be the
integration variable and set $z=(1-\xi)^2$. Sum over $i\ge1$ and
apply Jensen's inequality under $\sigma$ to each term of the mixture.
Subtracting the $i=1$ term gives
\begin{equation}\label{eq:low-marker-damping}
 \sum_{i\ge2}m_i\Delta_i
 \le e\sum_{i\ge2}i m_i(1-\sqrt e)^{i-1}.
\end{equation}
The sums are finite or positively dominated geometric sums, so
Tonelli justifies both operations.  Since $\rho_i\le1-e$ for
$i\ge2$, this bounds the negative part of the linear terms with $i\ge2$ in
\eqref{eq:low-master}.

For $q\le3$, discard $(1-s)^{i-2}\le1$ in
\eqref{eq:low-marker-damping}.  Bernoulli and monotonicity of
$y/(y-d)^{p+2}$ give
\[
 w_i\beta_i^p\le
        \frac{q^{p+2}(1+ci)}{\lambda(q+i)^{p+2}}.
\]
Expand $i=(q+i)-q$ and use~\eqref{eq:low-power-enclosures}, with
upper bounds on positive coefficients and lower bounds on negative
coefficients.  For $p=2$ this gives
$u^{-2}\sum_{i\ge2}i w_i\beta_i^2\le K_4$.
Thus the sum of the negative parts of $\Delta_i\rho_i S_i/u^2$ for $i\ge2$ is at most
$\tfrac32e(1-e)(1-s)K_4$.

For $3\le q\le10$, put $t_*=qs/(1-s)\le qd/(q+h)\le1$,
using~\eqref{eq:low-e-caps}.  Bernoulli and the convex chord
\[
 (1-t)+\frac{t}{1+y}-\frac1{1+ty}
       =\frac{t(1-t)y^2}{(1+y)(1+ty)}\ge0
\]
give
\[
 (1-s)^{i-1}\le\frac1{1-s}
       \left\{(1-t_*)+\frac{t_*q}{q+i}\right\}.
\]
Use the preceding Bernoulli weight bound with $p=4$.
Expanding $i$ and $i^2$ in powers of $q+i$, and applying
\eqref{eq:low-power-enclosures} term by term, gives the
bounds
\[
 \frac1{u^2}\sum_{i\ge2}i w_i\beta_i^4\le K_6,
 \quad\frac1{u^2}\sum_{i\ge2}\frac{qi w_i\beta_i^4}{q+i}\le K_7,
 \quad\frac1{u^2}\sum_{i\ge2}i^2w_i\beta_i^4\le K_{\rm sq}.
\]
Equation~\eqref{eq:low-marker-damping}, $\Delta_i\le i e$, and
\eqref{eq:low-all-age-prices} now bound the sums of the negative parts
of $\Delta_i\rho_i S_i/u^2$ and $\Delta_i^2J_i/u^2$, for $i\ge2$, by
$4e(1+s)\{(1-t_*)K_6+t_*K_7\}$ and $8e^2K_{\rm sq}$,
respectively.

It remains to bound the terms at $i=1$ and the term $C_{\rm cal}e$.
Combine the terms $F_{\rm pipe}+eS_1+e^2V_1$ before estimating:
\[
 F_{\rm pipe}+eS_1+e^2V_1
 =(1-e)^2V_1+u^2W_1(R)-2RuL_1+eB_1.
\]
The lower bound for the tail $V_1$ gives
$V_1/u^2\ge V_0+q\theta V_R$, and
$W_2^Y/(2u^2)=3k\lambda Rr_1/2\le3a_U Rr_1/2$.
To estimate $C_{\rm cal}$, replace $e$ by $E_v$ only in a
componentwise lower bound for the positive coefficient
$C_{\rm head}(2-e)$.  At $a_2=\beta_2(1-E_v)$,
\[
 A_0=E_v,\qquad
 2-(3\lambda+1)u+\lambda(\lambda+1)a_2=(2-h)E_v.
\]
The remaining factors are bounded below by
$u-\lambda\beta_2$ and
\[
 -k\log(\beta_2/u)-1+\beta_2/u\ge\phi>0.
\]
For the last inequality use
$-\log z\ge2(1-z)/(1+z)$ and $k\ge a_L/\lambda$, then insert
$\beta_2/u=(q+1)/(q+2+c)$.  This yields exactly the displayed
$\phi$. Its numerator is positive in both ranges of $q$:
for $q\le3$ this follows from $c\le1/6$, while for
$3\le q\le10$, $c=h/q$ gives
$c^2(q+2+c)\le31/216<6d$.
Finally, the identity
\[
 \frac{\kappa-1}{Z_p}=\frac{\theta(2-h)-1}{2q}>0
\]
gives $C_{\rm cal}/u^2\ge C_0\{\theta(2-h)-1\}$.
Thus~\eqref{eq:low-master} proves
\[
 F_{\rm floor}(v)/u^2>\mathcal G_1(\theta)\quad(q\le3),
 \qquad
 F_{\rm floor}(v)/u^2>\mathcal G_2(\theta)\quad(3\le q\le10).
\]
Both $\mathcal G_1$ and $\mathcal G_2$ are affine in $\theta$, so
part~(iv) of Lemma~\ref{lem:low-Q12-arithmetic} gives positivity
throughout $[\theta_{\rm lo},\theta_{\rm up}]$ in the respective
ranges of $q$. Since $1/15>1/29$, this covers $1/15<q\le10$
when $0<\alpha<1$.

\subsubsection{\texorpdfstring{Bounds for $F_{\rm pipe}$, $S_i$, and $J_i$ when $q\ge10$}{Bounds for F\_pipe, S(i), and J(i) when q >= 10}}
For $q\ge10$ we use the simpler bound
\begin{equation}\label{eq:low-large-master}
 F_{\rm floor}(v)>F_{\rm pipe}-eH_{\rm head}-e^2J_{\rm minus},
 \quad
 H_{\rm head}=\sum_{i\ge1}i[-S_i]_+,\quad
 J_{\rm minus}=\sum_{i\ge2}i^2[-J_i]_+.
\end{equation}
This follows from~\eqref{eq:low-master}, $\rho_i\le1$,
$\Delta_i\le ie$, and positivity of the coefficient of $e^2$ in braces
and of $C_{\rm cal}e$.
The next estimates will in particular show that the two negative-part
sums have finite support.

Put $x=1/q$, $c=hx$, $\lambda=1+hx$ and $g=q-R$.
We have $0<g<1$. Indeed, summing the lower
bound $z/(e^z-1)\ge1-z/2$ over the geometric series gives
\[
 T_p\ge\frac{h}{\lambda\log\lambda}-\frac1{2\lambda}
       >\frac{2q}{\lambda+1}-\frac1{2\lambda}.
\]
Here $k>2/(\lambda+1)$ follows by strict trapezoid comparison for
$\int_1^\lambda t^{-1}\dd t$.
Subtract the value $2q(q-1)/[(\lambda+1)q-\lambda]$ corresponding
to $R=q-1$.  The difference is
\[
 \frac{\lambda(\lambda+1)-qc^2}
 {2\lambda(\lambda+1)((\lambda+1)q-\lambda)}>0,
\]
since $qc^2=h^2/q\le1/40$.  The quotient is increasing in $T_p$,
so $R>q-1$; the upper bound $R<q$ was already proved.

We express $F_{\rm pipe}$ in terms of $T_p$ and its derivative,
\begin{equation}\label{eq:low-whole-source}
 F_{\rm pipe}=\frac{2q}{Z_p}\{-cT_p'-2T_p+T_p^2/q\},
 \qquad
 T_p=\frac{h^2}\lambda
                  \sum_{N\ge1}\frac{N d^{N-1}}{\lambda^N-1}.
\end{equation}
The derivative holds $h,d$ fixed.  To verify the identity, sum
$W_i^0$ to obtain $F_0=-cT_p'-(1+c/\lambda)T_p$.
Shifting the $W_i^p$ terms, retaining their first terms, gives
\[
 \sum_iw_iW_i^p=(\lambda+1)F_0+2T_p-2u^2-2(q-u)/\lambda.
\]
Insert this into the definition of $F_{\rm pipe}$ and use
$L_1=1/\lambda-u$.  The result is~\eqref{eq:low-whole-source}.
The same geometric majorants used in Lemma~\ref{lem:low-Lambert-tails}
justify the derivative and the summations.

Define
\begin{equation}\label{eq:low-source-rationals}
 \begin{gathered}
 f(h)=\frac{(1-h)(2h+1)}6,\quad
 s(a)=a-x/2+\frac{x^2(2-h)}{12a},\\
 s_l(a)=s(a)-\frac{x^4(24-36h+14h^2-h^3)}{720a^3},\\
 \begin{aligned}
 P_M(a)={}&2\{s(a)^2-(\lambda+1)s(a)+a^2-x^2(2-h)/12\}\\
 &-Mx\lambda\{2\lambda-(\lambda+1)s_l(a)\}.
 \end{aligned}
 \end{gathered}
\end{equation}
Applying~\eqref{eq:low-global-F}--\eqref{eq:low-global-C}
for each $N\ge1$ gives
\[
 \frac{s_l(a)}\lambda\le\frac{T_p}{q}\le\frac{s(a)}\lambda,
 \qquad
 K_{\rm actual}:=\frac{ch^2}{\lambda^2\delta^2}
       \sum_{N\ge1}d^{N-1}
         \frac{(N\delta)^2}{2(\cosh(N\delta)-1)}
 \ge\frac{a^2}{x\lambda^2}-\frac{x(2-h)}{12\lambda^2}.
\]
Writing $t_T=T_p/q$, equation~\eqref{eq:low-whole-source} becomes
\[
 F_{\rm pipe}=
 \frac{2\{t_T^2/x-(1+1/\lambda)t_T/x+K_{\rm actual}\}}
                 {2-(\lambda+1)t_T}.
\]
Since $s(a)\le1-x/4+x^2/6<1$, the quadratic in $t_T$ decreases
on the whole required interval.  Its lower bound uses $s(a)/\lambda$,
and its denominator is at most $2-(\lambda+1)s_l(a)/\lambda$.
For either positive target $M=f(h)-2x/5$ or $M=69/500$,
$P_M(a)>0$ therefore implies $F_{\rm pipe}>M$: the inequality
first proves positivity of the lower numerator, so replacing the
denominator preserves the lower bound. For $a_L\le a\le a_U$,
$s\ge19/20$, $\lambda+1\le41/20$, $0<s'<1$, and $s_l'>0$.
Consequently $P_M'(a)>37/10$.  It suffices to check $P_M(a_L)>0$.

To bound $S_i=B_i-2V_i$, set $t=i/q$, let
$z=d\lambda^{-i}$, and put $w=hz/(1-z)$.  Bernoulli gives
$0<w\le d/(1+t)$.  For a free point in this larger domain define
\begin{align}
 b&=xw,\quad r=(\lambda+b)^{-1},\quad r_{\rm prev}=(1-b)/\lambda,
 \nonumber\\
 R&=1/x-g,\qquad
 \alpha_p=\frac{\lambda^2+gd(\lambda+1)x^2}{1+x},
 \nonumber\\
 A&=1-b+(cb+b^2)/2,\nonumber\\
 P&=c+2b-4cb-5b^2+\tfrac92c^2b+\tfrac{25}2cb^2+8b^3,\nonumber\\
 C&=5b/3+cz/3,\nonumber\\
 U&=\{i(A+RP)+z/2+RC\}/\lambda,\nonumber\\
 \mathcal B&=2i(\alpha_p+cR)/\lambda+2iRrb/\lambda
             +2(i+1)Rr^2b\nonumber\\
 &\quad+\frac2{\lambda+1}(2i-1)
                \{1-(\alpha_p+\lambda R)r+(\lambda+1)R-R/r_{\rm prev}\},
 \nonumber\\
 \operatorname{shape}&=2ht+4(t-1/3)w-(2h+3)z/3,
 \qquad z=w/(h+w).\label{eq:low-large-price-rationals}
\end{align}
At $t=i/q$, $z=d\lambda^{-i}$ and $w=hz/(1-z)$, we have
$B_i/w_i\ge\mathcal B$ and $V_i/w_i\le U$.
The first follows from $k\ge2/(\lambda+1)$ and $\tau_i\ge0$.
For the second bound, the coefficients in
\eqref{eq:low-positive-mixture} satisfy, for every integer $N\ge1$,
\begin{align*}
 A_N&\le1-(N-1)c/2+(N^2-1)c^2/12,\\
 P_N&\le Nc-(N-1)(5N+2)c^2/6+(N^2-1)(4N+1)c^3/12,\\
 B_N&\le(N-1)/(2N),\qquad
 C_N\le c(N-1)(5N+2)/(6N).
\end{align*}
For the first bound, multiplication by $(1+c)^N-1$ leaves, at
power $c^j$, $j\ge4$, the nonnegative coefficient
\[
 \binom{N}{j-2}
 \frac{(j-3)(N+1)\{(j-4)N+5j-8\}}{12j(j-1)};
\]
the lower powers cancel.  The second bound follows by inserting the
first in $P_N$ and using the global third-order Taylor lower bound
for $(1+c)^{1-N}$, whose fourth derivative is nonnegative.
For the third, put
$D_N=(N-1)Q_N^2-2N\sum_{j=1}^{N-1}j\lambda^j$.
One has $D_N(1)=0$,
$D_N'(1)=N^2(N-1)(N-2)/3\ge0$.
The coefficients of $D_N'$ are negative only at degrees at most
$N-2$ and positive above that degree: the coefficient at degree
$j-1$ is
$j\{(N-1)\min(j+1,2N-1-j)-2Nj\mathbf1_{j<N}\}$.
Thus $D_N'/\lambda^{N-2}$ increases on $[1,\infty)$ and
$D_N\ge0$, proving the $B_N$ bound.
For the last bound we use the polynomial identity
\begin{align*}
 &(\lambda+1)N^2\lambda^{2N-1}-cN\lambda^NQ_N-(2+2Nc)Q_N^2
       =c^2\sum_{j=0}^{2N-2}a_j\lambda^j,\\
 &a_j=(j+1)(j+2)(N-1-j/3)\quad(0\le j\le N-2),\\
 &a_{2N-2-r}=(N-r-1)\{(2r+1)N+r(r+1)\}+2r(r+1)(r+2)/3,\\
 &\hspace{48mm}0\le r\le N-1.
\end{align*}
For $N=1$ both sides vanish.  For $N\ge2$ all $a_j\ge0$;
finite second differences, with the four joining coefficients,
verify the identity.  Their sum is
$N^2(N-1)(5N+2)/6$.  Since
$C_N=c\lambda^{1-N}\sum_j a_j\lambda^j/(NQ_N^2)$ and
$Q_N\ge N\lambda^{(N-1)/2}$, the required upper bound follows.
Summing all four inequalities under $\Pi$ gives exactly the
$A,P,z/2,C$ in~\eqref{eq:low-large-price-rationals} for every $i\ge1$.

To estimate the negative parts of $J_i$, we use one further rational function. For
$q\ge10$, $0<h\le1/2$ define
\begin{equation}\label{eq:low-AJ-rational}
 \begin{gathered}
 a_m=\frac{d}{\lambda(q+2)},\qquad
 r_m=\frac1{\lambda+d/(q+2)},\\
 N_m=\frac{(\lambda a_m+c)
       \{2\lambda^2a_m^2+(3\lambda^2+2\lambda+1)a_m+\lambda^2-1\}}
                 {\lambda^2(1-\lambda a_m)},\\
 A_J^{\rm lo}=\lambda r_m^2+
                   \frac{4\lambda^2u r_m}{\lambda+1}-qN_m.
 \end{gathered}
\end{equation}

\begin{lemma}[Rational inequalities for $q\ge10$]
\label{lem:low-Q3-arithmetic}
For $0<x\le1/10$, $0\le h\le1/2$,
$P_{f(h)-2x/5}(a_L)>0$.
For the domain
\[
 0<x\le1/10,\quad0<h\le1/2,\quad t\ge0,\quad
 0<w\le(1-h)/(1+t),\quad0\le g\le1,
\]
the rational functions in~\eqref{eq:low-large-price-rationals} satisfy
\begin{equation}\label{eq:low-Q3-price-check}
 \mathcal B-2U\ge\operatorname{shape}-\frac{28}{5}x(1+ht).
\end{equation}
Finally, $A_J^{\rm lo}>3-12/q$ for $q,h>0$ with $h\le1/2$.
\end{lemma}

These checks have the following finite specifications.
Put $A=12a_L$.  The polynomial
$1800A^3P_{f(h)-2x/5}(a_L)/x^3$ has constant coefficient $380160$;
the sum of the absolute values of its negative coefficients, weighted
by $10^{-j}2^{-k}$ for $x^jh^k$, is
$65345524702337817/204800000000<320000$.
For~\eqref{eq:low-Q3-price-check}, multiply
$[\mathcal B-2U-\operatorname{shape}]/x+28(1+ht)/5$ by the
positive denominator
\[
 3(h+w)(1+x)\lambda(\lambda+1)(1-xw)(\lambda+xw)^2.
\]
Substitute $w=(1-h)v_0/(1+t)$ and multiply by $(1+t)^7$.
The result is a polynomial of degrees $(6,8,8,7)$ in
$(x,h,t,v_0)$, affine in $g$.  At $g=0,1$, transform only
$x,h,v_0$ by~\eqref{eq:low-Bernstein-transform} on
$[0,1/10]\times[0,1/2]\times[0,1]$, keeping ordinary powers of
$t\ge0$.  Each endpoint has $4319$ nonzero positive coefficients,
with minimum $69/280$; the remaining coefficients are zero.
For~\eqref{eq:low-AJ-rational}, the numerator of
$A_J^{\rm lo}-(3-12/q)$ over
\[
 q(q+h)^3(2q+h)(q+1)(q+2)^2(q+h+1)(q^2+3q+2h)^2
\]
has $60$ nonnegative nonzero coefficients in $q,h$, including
$66q^{11}$.  These specifications prove the three statements by
finite rational arithmetic, with every denominator positive.

The first assertion gives
\begin{equation}\label{eq:low-Q3-source}
 F_{\rm pipe}>f(h)-2/(5q).
\end{equation}
To apply~\eqref{eq:low-Q3-price-check} to $S_i/w_i$, we first
show that for $t=i/q$, $z=d\lambda^{-i}$ and $w=hz/(1-z)$,
\begin{equation}\label{eq:low-Q3-shape}
 \operatorname{shape}\ge\frac{9t-7}{3(1+t)}+\frac34ht.
\end{equation}
For $t\le1/3$, use $w\le d/(1+t)$ and $z\le d$; the difference
in~\eqref{eq:low-Q3-shape} is at least
$h\{5t/4-4(t-1/3)/(1+t)+1/3\}+2h^2/3\ge0$.
For $t\ge1/3$ and $y=ht<3$, the exponential series, using
$n!\ge2\cdot3^{n-2}$ for $n\ge2$, gives
$e^y\le1+y+y^2/[2(1-y/3)]$.
Since $\lambda^i\le e^y$ and $\lambda^i\ge1+y$,
\[
 w\ge\frac{d(1-y/3)}{1+t-y/3+yt/6},\qquad z\le\frac d{1+y}.
\]
After insertion, the difference is at least
\[\begin{aligned}
 &\frac{h}{12(t+1)(y+1)((y+6)t+6-2y)}\\
 &\quad\times\Bigl\{t^3(15y^2-27y+18)+t^2(81y^2+13y+108)\\
 &\hspace{19mm}+t(42y^2-48y+18)+120+152y-40y^2
                  +\frac{48y-16y^2}{t}\Bigr\}>0.
\end{aligned}\]
The first and third quadratic factors have negative discriminants,
the second has positive coefficients, the constant quadratic has
positive endpoint values on $[0,3]$, and the last term is
nonnegative.  The case $y=3$ follows by continuity.
For $y\ge3$, $t\ge6$; drop the nonnegative $w$ term and use
$(2h+3)d\le3$ to get
$\operatorname{shape}\ge2y-1/(1+y)>3+3y/4$.
This proves~\eqref{eq:low-Q3-shape} in every case.  Thus, with
$\eta=28/(5q)$,
\begin{equation}\label{eq:low-Q3-head-envelope}
 \frac{S_i}{w_i}\ge\frac{9t-7}{3(1+t)}-\eta
                                  +ht(3/4-\eta).
\end{equation}
Since $\eta\le14/25<3/4$, $S_i<0$ is possible only if
$t<(7+3\eta)/(9-3\eta)$.

For $i\ge2$, the lower bound from the mixture and the expression
for $J_i$ give
\[
 J_i/w_i>iA_J-B_J,
 \quad A_J=\lambda r_i^2+2k\alpha_pr_i
 +R\{cr_i^2+2\lambda r_i^3L_i+k(\lambda r_i-1/r_{i-1})\},
\]
where $B_J=k\{C_\alpha r_i-R(1/r_{i-1}-\lambda r_i)/2\}<q+1$.
To bound $A_J$, set $a=L_i$.  At $k=1$ its bracket is exactly
\[
 -\frac{(\lambda a+c)
       \{2\lambda^2a^2+(3\lambda^2+2\lambda+1)a+\lambda^2-1\}}
       {\lambda^2(1+a)^3(1-\lambda a)}.
\]
Since $k\le1$, this is a lower bound for that bracket.
Bernoulli gives $a\le a_m$, $r_i\ge r_m$ for $i\ge2$.
Increase every positive numerator factor to $a_m$, omit the factor
$(1+a)^3\ge1$, and use $k\ge2/(\lambda+1)$,
$\alpha_p>\lambda^2u$, and $R<q$.  The result is
$A_J\ge A_J^{\rm lo}>3-12/q$ by Lemma~\ref{lem:low-Q3-arithmetic}.
Also $w_i\le w_1=u^2<1$, since
$y/(1-dy)^2$ increases in $y=\lambda^{-i}$.
Hence a finite sum bounds $J_{\rm minus}$. For $T\ge2$, writing
$n=\lfloor T\rfloor$ and $r=T-n$,
Faulhaber's formulas give
\begin{align*}
 \sum_{i=1}^n i^2(T-i)
 ={}&T^4/12+T^2(-r^2/2+r/2-1/12)\\
 &+T(2r^3/3-r^2+r/3)-r^2(1-r)^2/4
 \le(T^4+T^2)/12.
\end{align*}
The coefficient of $T^2$ is at most $1/24$, that of $T$ is at most
$1/12$ in absolute value, and $T/12\le T^2/24$ for $T\ge2$.
Taking $T=(q+1)/(3-12/q)$ gives
\begin{equation}\label{eq:low-Q3-square-cost}
 J_{\rm minus}<q^4J(q),\qquad
 J(q)=\frac{(q+1)^4}{324q(q-4)^3}
                       +\frac{(q+1)^2}{36q^3(q-4)}.
\end{equation}
For every integer $i>T$ the bound on $J_i/w_i$ gives $J_i>0$.

\subsubsection{Bounds for \texorpdfstring{$q\ge21$ and for $q\ge10$ with $h\ge1/4$}{q >= 21 and for q >= 10 with h >= 1/4}}
We first give a quadrature bound that includes a support endpoint
lying inside a mesh interval.  If $G$ is continuous, compactly
supported in $[0,\infty)$, $G(0)=0$, and $G'$ has bounded variation
on $(0,\infty)$, then
\begin{equation}\label{eq:low-BV-quadrature}
 \left|\Delta\sum_{i\ge1}G(i\Delta)-\int_0^\infty G(t)\dd t\right|
 \le\frac{\Delta^2}{8}\operatorname{TV}(G').
\end{equation}
On a mesh interval $[a,b]$ use $K(s)=(s-a)(b-s)/2$ and integrate
twice by parts:
\[
 \frac{\Delta}{2}\{G(a)+G(b)\}-\int_a^bG(s)\dd s
                  =\int_{(a,b)}K(s)\dd G'(s).
\]
Derivative jumps contribute their point masses; $K$ vanishes at
both mesh endpoints.  Since $0\le K\le\Delta^2/8$, summing proves
\eqref{eq:low-BV-quadrature}.  The first and last values of $G$ are
zero, so the complete trapezoid sum equals the right-endpoint sum.

For $t=i/q$, $i\ge1$, monotonicity of $s/(s-d)^2$ and
$\lambda^i\ge1+ht$ give
\begin{equation}\label{eq:low-Bernoulli-weight}
 w_i\le\frac{1+ht}{\lambda(1+t)^2}.
\end{equation}
If $q\ge21$, use $\eta_A=4/15$, $A=39/5$, $B=41/5$ in
\eqref{eq:low-Q3-head-envelope}. Every index for which $S_i<0$ satisfies
$t<A/B=39/41<1$.
Put $D_h=1-(\lambda+1)h$.
On that support $D_h(1+ht)\le(1-2h)(1+ht)\le1$.
The bound on $eq^2$ in~\eqref{eq:low-e-caps} and
\eqref{eq:low-Bernoulli-weight} imply
\[
 eH_{\rm head}\le\frac1q\sum_{i\ge1}G_A(i/q),\qquad
 G_A(t)=\left[\frac{t(A-Bt)}{3(1+t)^3}\right]_+.
\]
Inside its support,
\[
 G_A'=\frac{A-32t+Bt^2}{3(1+t)^4},\qquad
 G_A''=\frac{-32-4A+(96+2B)t-2Bt^2}{3(1+t)^5}.
\]
The numerator of $G_A''$ increases across the support, is negative
at zero and positive at its end.  Thus $G_A'$ has one minimum,
then jumps from a negative value to zero at the end.
Moreover $G_A'>-1/2$: after multiplying by $6(1+t)^4$ the
required inequality is
\[
 17-52t+36t^2+12t^3+3t^4+6\eta_A(1-t^2)>0.
\]
For $t\le1/3$ its first three terms are at least $11/3$; for
$t\ge1/3$, use $12t^3\ge4t^2$ and
$17-52t+40t^2>0$, whose discriminant is $-16$.
The last group is nonnegative because $t<1$.
Consequently $\operatorname{TV}(G_A')<13/5+1<4$.
An elementary antiderivative gives
\[
 I_A:=\int G_A=\frac83-\frac{B^2}{96}-\frac B3\log(16/B).
\]
For $0<s<1$ the exact geometric remainder in
$\int_0^s(1-r^2)^{-1}\dd r$ gives
\begin{equation}\label{eq:low-log-four}
 \log\frac{1+s}{1-s}>\ell(s):=2(s+s^3/3+s^5/5+s^7/7).
\end{equation}
Thus, with
\[
 I_A^{\rm up}=175/96+(9/16)\eta_A-(3/32)\eta_A^2
                    -(3-\eta_A)\ell(39/121),
\]
equation~\eqref{eq:low-BV-quadrature} gives
$eH_{\rm head}<I_A^{\rm up}+1/(2q^2)$.
The same bound for $J_{\rm minus}$, with $3-12/q\ge17/7$, yields
\[
 e^2J_{\rm minus}<
 J_A:=\frac{(22/21)^4}{12(17/7)^3}
               +\frac{(22/21)^2}{12(17/7)21^2}
       =\frac{2904605}{409493637},
\]
because $e^2q^4<1$.
Combining these bounds with $f(h)\ge1/6$ in~\eqref{eq:low-Q3-source}
gives
\[
 F_{\rm floor}(v)>
 \frac16-\frac2{105}-I_A^{\rm up}-\frac1{882}-J_A
 =\frac{28155382207511429143103}{124404112403316879470013600}
 >\frac1{5000}.
\]

Next let $q\ge10$ and $h\ge1/4$.  Use $\eta_B=14/25$,
$A=217/25$, $B=183/25$. When $S_i<0$ we have $t<A/B<2$.
Retain the factor $D_h$ in the bound on $eq^2$ in~\eqref{eq:low-e-caps}. Equations
\eqref{eq:low-Q3-head-envelope} and~\eqref{eq:low-Bernoulli-weight} give
\[
 eH_{\rm head}\le\frac{D_h}{q}\sum_{i\ge1}G_B(i/q),\qquad
 G_B(t)=\left[\frac{t(A-Bt)(1+t/2)}{3(1+t)^3}\right]_+.
\]
Inside the support,
\[\begin{aligned}
 G_B'&=\frac{2A-(2A+4B)t-(A+B)t^2}{6(1+t)^4},\\
 G_B''&=\frac{-10A-4B+(4A+10B)t+2(A+B)t^2}{6(1+t)^5}.
\end{aligned}\]
The second derivative changes sign once from negative to positive
before $t=1$, so $G_B'$ again has one minimum and its last jump is
to zero.  The inequality $G_B'>-1$ reduces to
\[
 584/25-(566/25)t+20t^2+24t^3+6t^4>0;
\]
the quadratic part has discriminant $-847644/625<0$.
Therefore $\operatorname{TV}(G_B')<217/75+2<5$.
The exact integral is $I_B=(8/3)\log(400/183)-73563/40000$.
For $s_B=217/583$ the remainder of the geometric series gives
\[
 \log(400/183)<2\{s_B+s_B^3/3+s_B^5/[5(1-s_B^2)]\}.
\]
Together with~\eqref{eq:low-BV-quadrature} and $q\ge10$, this proves
\[
 eH_{\rm head}<D_h C_B,\qquad
 C_B=\frac{122243248244149}{483498900280000}>\frac1{12}.
\]
The bound for $J_{\rm minus}$ with $3-12/q\ge9/5$ gives
\[
 e^2J_{\rm minus}<D_h^2J_B,\qquad
 J_B=\frac{(11/10)^4}{12(9/5)^3}
       +\frac{(11/10)^2}{12(9/5)10^2}
       =\frac{187913}{8748000}.
\]
Since $D_h\le1-2h$, equations~\eqref{eq:low-large-master} and
\eqref{eq:low-Q3-source} imply
\[
 F_{\rm floor}(v)>f(h)-\frac1{25}-(1-2h)C_B-(1-2h)^2J_B.
\]
The right side has derivative
$(1-4h)/6+2C_B+4(1-2h)J_B>-1/6+2C_B>0$ on $[1/4,1/2]$.
At $h=1/4$ it equals
\[
 \frac{16616808991576441}{1057412094912360000}>\frac1{100}.
\]
This proves positivity for $q\ge10$ and $h\ge1/4$, using the same
$e$ in the bounds for $eH_{\rm head}$ and $e^2J_{\rm minus}$.

\subsubsection{\texorpdfstring{The coefficients $S_1,\ldots,S_{24}$ for $10\le q\le21$ and $h\le1/4$}{The first 24 coefficients S(i) for 10 <= q <= 21 and h <= 1/4}}
Put $x=1/q\in[1/21,1/10]$. Since $25/21>217/183$,
\eqref{eq:low-Q3-head-envelope} gives $S_i\ge0$ for $i\ge25$.
It remains to bound $S_i$ for $1\le i\le24$.

We first bound $g$ and $F_{\rm pipe}$. With
$s,s_l,P_M$ as in~\eqref{eq:low-source-rationals}, let
\[
 t_g(\gamma)=\frac{2(1-\gamma x)}{x+(\lambda+1)(1-\gamma x)}.
\]
For $0<x\le1/10$, $0\le h\le1/4$, the finite inequalities
needed are
\begin{equation}\label{eq:low-Q3c-source-checks}
 s_l(a_L)/\lambda>t_g(1/4),\qquad
 t_g(1/8)>s(a_U)/\lambda,\qquad P_{69/500}(a_L)>0.
\end{equation}
They imply $1/8<g<1/4$ and $F_{\rm pipe}>69/500$ by monotonicity
of $R$ in $T_p$ and the bound using $P_M$ preceding
Lemma~\ref{lem:low-Q3-arithmetic}.

We next define a polynomial for each of the remaining indices.
To form $F_i(c,x,g)$ take $\lambda=1+c$, $d=1-c/x$ and
$a_U=1+c/2-c^2/12+c^3/24$; substitute $c=hx$ at the end.  For each
integer $0\le n\le25$ put
\begin{equation}\label{eq:low-Q3c-Dn}
 D_n(c,x)=1+x\sum_{j=1}^n\binom nj c^{j-1}.
\end{equation}
The empty sum is zero.  When $c=hx$ this is
$(\lambda^n-d)/h=\beta_n^{-1}$, with its continuous value at
$h=0$.  For $1\le i\le24$ let $n=i+1$, $p=\lambda^n$,
$b_0=1/D_n$, and set
\begin{align*}
 m_0&=b_0,\qquad m_1=pb_0^2,\qquad m_2=p(p+d)b_0^3,\\
 m_3&=p(p^2+4dp+d^2)b_0^4,\\
 m_4&=p(p^3+11dp^2+11d^2p+d^3)b_0^5,\\
 r&=D_i/D_{i+1},\qquad r_{\rm prev}=D_{i-1}/D_i,\\
 L_i&=dx/(\lambda D_i),\qquad w_i=\lambda^{i-1}/D_i^2.
\end{align*}
For a free $a\ge1$, use the same $R=1/x-g$ and $\alpha_p$ as in
\eqref{eq:low-large-price-rationals} and define
\begin{align}
 \operatorname{TL}_i(a)&=
       \{am_0/x+m_1/2+x m_2/(12a)-x^3m_4/(720a^3)\}/\lambda,
 \nonumber\\
 \operatorname{PU}_i(a)&=\lambda^{-2}
 \Bigl\{n(am_1+x m_2/2+x^2m_3/(12a))+a^2m_0/x\nonumber\\
 &\hspace{30mm}-xm_2/12+x^3m_4/(240a^2)\Bigr\},
 \nonumber\\
 \operatorname{VU}_i(a)&=-(1+cR)\operatorname{TL}_i(a)
              +\{1+(\lambda+1)R\}\operatorname{PU}_i(a)\nonumber\\
 &\quad-2Rn\lambda^i/D_{i+1}^2-2R/(x\lambda D_{i+1}),
 \nonumber\\
 \operatorname{BB}_i(a)&=2i(\alpha_p+cR)/\lambda+2iRrL_i
       +2nR\lambda r^2L_i\nonumber\\
 &\quad+(a/\lambda)(2i-1)
              \{1-(\alpha_p+\lambda R)r+(\lambda+1)R-R/r_{\rm prev}\},
 \nonumber\\
 K(t,x)&=\frac{(3-9x)t-7/3+3x}{1+t},\nonumber\\
 F_i(c,x,g)&=\operatorname{BB}_i(a_U)
                   -2\operatorname{VU}_i(a_U)/w_i-K(ix,x).
 \label{eq:low-Q3c-head-rational}
\end{align}
Each $F_i$ is an explicit rational function. To clear its denominators, set
\begin{equation}\label{eq:low-Q3c-head-denominator}
 \mathcal D_i=720\lambda^{i+1}x^2a_U^3D_{i+1}^5
                                (1+x)D_{i-1}(1+ix).
\end{equation}
For $\gamma\in\{1/8,1/4\}$ define the polynomial
\begin{equation}\label{eq:low-Q3c-head-polynomial}
 \mathcal P_{i,\gamma}(x,h)
 =x^{-2}\left.\{\mathcal D_i F_i(c,x,\gamma)\}\right|_{c=hx}.
\end{equation}
For each $1\le i\le24$, the arithmetic lemma checks that the numerator
formed using~\eqref{eq:low-Q3c-Dn} is divisible by $x^2$.

\begin{lemma}[Rational inequalities for $0<x\le1/10$ and $0\le h\le1/4$]
\label{lem:low-Q3c-arithmetic}
The three inequalities~\eqref{eq:low-Q3c-source-checks} hold on
$0<x\le1/10$, $0\le h\le1/4$.
Each of the $48$ polynomials~\eqref{eq:low-Q3c-head-polynomial}
has nonnegative tensor Bernstein coefficients on
$[0,1/10]\times[0,1/4]$, and is strictly positive when $x>0$.
Consequently $F_i(hx,x,g)>0$ for $1\le i\le24$,
$0<x\le1/10$, $0\le h\le1/4$, and $1/8\le g\le1/4$.
\end{lemma}

To specify the three smaller checks by positive denominators, put
$A=12a_L$, $A_u=24a_U$,
$F_4=4x+(\lambda+1)(4-x)$ and
$F_8=8x+(\lambda+1)(8-x)$.
The cleared polynomials are
\[
 \frac{60\lambda A^3F_4}{x^2}\{s_l(a_L)/\lambda-t_g(1/4)\},\quad
 \frac{24\lambda A_uF_8}{x^2}\{t_g(1/8)-s(a_U)/\lambda\},\quad
 \frac{90000A^3}{x^2}P_{69/500}(a_L).
\]
Their degrees in $(x,h)$ are $(8,9)$, $(6,7)$, $(9,10)$; all
$90$, $56$, and $110$ Bernstein coefficients are positive.  Their
respective minima are $25920$, $672/5$, and $9325152/125$.
For $\mathcal P_{i,\gamma}$, the degrees are $(27,20)$ at $i=1$ and
$(7i+20,7i+12)$ at $2\le i\le24$.
The coefficient operation is~\eqref{eq:low-Bernstein-transform};
the $636112$ nonzero coefficients are positive. In each array the
only zero coefficient is the one at index $(0,0)$, which proves
strict positivity for $x>0$. Checking the
identities~\eqref{eq:low-Q3c-head-rational}--\eqref{eq:low-Q3c-head-polynomial}
and these finite arrays is the entire arithmetic obligation for this
lemma.  Positivity of $D_n$, $a_U$, $\lambda$, $x$ and $1+ix$
proves the sign of~\eqref{eq:low-Q3c-head-denominator}. The factor
$x^2$ is positive for $x>0$. Affinity of $F_i$ in $g$
then covers the whole interval between the two checked endpoints.

We now use the polynomial inequalities to bound $S_i/w_i$. At
$a=c/\log\lambda$, Lemma~\ref{lem:low-Lambert-tails} gives
$V_i\le\operatorname{VU}_i(a)$ and
$B_i/w_i=\operatorname{BB}_i(a)$.
We must change this one common $a$ in both terms together.
For $i\le24$, $c\le1/40$, and $n=i+1\le25$,
\[
 p=\lambda^n<e^{5/8}
 \le1+5/8+\frac{(5/8)^2}{2(1-5/24)}=\frac{569}{304}<2.
\]
The exponential estimate includes the full tail, by the factorial
bound used in~\eqref{eq:low-Q3-shape}.  Moreover
$p-1<2nc$, so $b_0>1/(1+2n/q)\ge1/6$, while $b_0\le1$.
The displayed moments then give $m_3/m_1\le13$ and
$m_4/m_0\le150$.
Differentiating only the explicit rational functions yields, for all
$a\ge1$,
\[
 \operatorname{TL}_i'(a)<\frac{101m_0}{100x\lambda},\qquad
 \operatorname{PU}_i'(a)>\frac{19m_0}{10x\lambda^2}.
\]
Indeed, the derivative of the first has numerator
$m_0/x-xm_2/(12a^2)+x^3m_4/(240a^4)$; in the second its
$n$ term is positive since $x^2m_3/(12a^2m_1)\le13/1200<1$.
The remaining fourth-order ratios are at most $1/16000$ and
$1/8000$, respectively.  As $R>q-1/4$, $cR\le1/4$ and
$\lambda\le41/40$, it follows that
\[
 \operatorname{VU}_i'(a)>
 \frac{m_0}{x\lambda^2}
 \left\{\frac{19}{10}(2q+1/2)
                 -\frac{101}{100}\frac{41}{40}\frac54\right\}
 >3q^2m_0.
\]
For the last inequality, subtract $3q(41/40)^2$ from the bracket;
the result is $(1037/1600)q-1101/3200\ge19639/3200>0$.
We proved $\tau_i<1$ and $w_i<1$.  Hence
\[
 \frac{\dd}{\dd a}
       \{\operatorname{BB}_i(a)-2\operatorname{VU}_i(a)/w_i\}
 <47-6q^2b_0/w_i\le47-q^2<0.
\]
Now $a\le a_U$ and Lemma~\ref{lem:low-Q3c-arithmetic} imply
\begin{equation}\label{eq:low-Q3c-actual-head}
 \begin{split}
 S_i/w_i&\ge\operatorname{BB}_i(a)-2\operatorname{VU}_i(a)/w_i\\
 &\ge\operatorname{BB}_i(a_U)-2\operatorname{VU}_i(a_U)/w_i
 >K(ix,x),\qquad1\le i\le24.
 \end{split}
\end{equation}
This bounds $S_i/w_i$ at every index where $S_i$ can be negative.
For $i\ge25$ we already have $S_i\ge0$ by~\eqref{eq:low-Q3-head-envelope}.

To bound $H_{\rm head}$, we sharpen the estimate for $w_i$. Let
$y=i\log\lambda$ and $\rho=(1-h/2)/a\le1$.  Directly,
\[
 w_i=\frac{h^2}{\lambda
       \{h\cosh(y/2)+(2-h)\sinh(y/2)\}^2}
 \le\frac1{\lambda(1+\rho i/q)^2}.
\]
Also
\[
 \lambda(1-h/2)^2-d(1+c/2)^2
                  =\frac{h^2}{4}(\lambda-dx^2)>0.
\]
Since $a\le1+c/2$, this proves $\lambda\rho^2>d$ and
$\lambda\rho>d$.  All coefficients of
$\lambda(1+\rho t)^2-d(1+t)^2$ are therefore positive, so
\begin{equation}\label{eq:low-Q3c-weight}
 d w_i\le(1+i/q)^{-2}\qquad(i\ge1).
\end{equation}
Apply $e\le d/(q+1)^2$ to~\eqref{eq:low-Q3c-actual-head}.
Since $S_i\ge0$ for $i\ge25$, this gives
\[
 eH_{\rm head}\le\frac{u^2}{q}\sum_{i\ge1}G_x(i/q),\qquad
 G_x(t)=\left[\frac{t\{(7-9x)-(9-27x)t\}}{3(1+t)^3}\right]_+.
\]
For $0\le x\le1/10$ the support ends at
$T=(7-9x)/(9-27x)<1$.  With $A=7-9x$, $B=9-27x$,
\[
 G_x'=\frac{A-2(A+B)t+Bt^2}{3(1+t)^4},\qquad
 G_x''=\frac{-6A-2B+(6A+8B)t-2Bt^2}{3(1+t)^5}.
\]
The second derivative changes sign exactly once on the support.
Furthermore $G_x'>-1/2$.  The required numerator is
\[
 17-52t+36t^2+12t^3+3t^4+18x(-1+8t-3t^2).
\]
For $t\le1/4$, the quadratic part is at least $25/4$ and the
last term at least $-9/5$.  For $1/4\le t\le1$, the last
term is nonnegative; the same split at $1/3$ used for $G_A$
proves positivity.  Including the jump at $T$ gives
$\operatorname{TV}(G_x')\le7/3+1=10/3$.
Thus~\eqref{eq:low-BV-quadrature} gives
\[
 q^{-1}\sum_{i\ge1}G_x(i/q)\le I(x)+1/240,\qquad
 I(x)=\frac{A(A+2B)}{6(A+B)}-\frac B3\log\frac{A+B}{B}.
\]
For each fixed $t$, $G_x(t)$ is the positive part of an affine
function of $x$, so its integral is convex on $[0,1/10]$.
Its maximum is bounded by its two endpoint values.  Use
\eqref{eq:low-log-four} at $s=7/25$ and $61/187$ to obtain
\begin{align*}
 I(0)&<175/96-3\ell(7/25)=\frac{56737364551}{585937500000},\\
 I(1/10)&<11407/7440-(21/10)\ell(61/187)
              =\frac{33081533096174793233}{297463843858303887600}.
\end{align*}
The second rational bound exceeds the first by
\[
 \frac{696221180224885514072002509}{48415339169645815039062500000}>0.
\]
Consequently
\[
 eH_{\rm head}<u^2C_H,\qquad
 C_H=\frac{17160482889458863049}{148731921929151943800}.
\]
The same bound on $e$ and~\eqref{eq:low-Q3-square-cost} give
$e^2J_{\rm minus}<d^2u^4J(q)\le u^4J(q)$.
Together with $F_{\rm pipe}>69/500$, equation~\eqref{eq:low-large-master}
therefore proves
\[
 F_{\rm floor}(v)>L(q):=69/500-u^2C_H-u^4J(q).
\]
For the last continuous comparison, ordinary differentiation yields
\[
 \frac{(q+1)^3}{q}L'(q)
 =-2C_H+\frac{q(q+1)^3}{27(q-4)^4}
              +\frac{q^2-2q+2}{18q(q-4)^2}.
\]
Both nonconstant terms decrease strictly for $q\ge10$.
For the first, its logarithmic derivative is
$1/q+3/(q+1)-4/(q-4)<0$; the derivative of the second is
$(-q^3-6q+8)/[18q^2(q-4)^3]<0$.
Thus $L'$ can change sign only from positive to negative, and
$L$ has no interior minimum on $[10,21]$.  Its two endpoint
margins are
\[\begin{gathered}
 L(10)-1/40=\frac{195098675960795601986623}{65597470507242819053271000}>0,
 \\
 L(21)-1/40=\frac{713805528405149594753}{359931251068547703996000}>0.
\end{gathered}\]
This proves $F_{\rm floor}(v)>1/40$ for $10\le q\le21$ and $h\le1/4$.

\begin{proposition}[Negativity of $v'(\lambda)$ for $1<\lambda<\min(m,2/\sqrt3)$]\label{prop:entire-low-region}
For every fixed offspring law satisfying the assumptions of Theorem~\ref{thm:computer},
$v'(\lambda)<0$ for $1<\lambda<\min(m,2/\sqrt3)$.
\end{proposition}
\begin{proof}
Lemma~\ref{lem:calibration} covers $\alpha=1$. When $0<\alpha<1$,
we have $q>0$ and $0<h<1/2$. The six ranges
\[\begin{gathered}
 q\le1/15;\qquad1/15<q\le3;\qquad3<q\le10;\qquad q\ge21;\\
 10<q<21,\ h\ge1/4;\qquad10<q<21,\ 0<h<1/4
\end{gathered}\]
exhaust these parameters.  The preceding arguments give
$F_{\rm floor}(v)>0$ in the last five, and directly $D_\lambda>0$
in the first.  Lemma~\ref{lem:mellin-floor} gives
$D_\lambda\ge F_*>F_{\rm floor}(v)$ in all remaining cases.
Since $Z>0$, the derivative formula~\eqref{eq:full-sign} gives
$v'(\lambda)<0$. The finite arithmetic inputs are
exactly Lemmas~\ref{lem:low-Q0-arithmetic},
\ref{lem:low-Q12-arithmetic}, \ref{lem:low-Q3-arithmetic}, and
\ref{lem:low-Q3c-arithmetic}; their expressions, domains, positive
denominators, and uses in the proof have all been specified above.
\end{proof}
